\begin{exsol}{17}{\MG{} Ch.~4 Exercise 12}{185--186}{Exam}{L3}
\asked With $p(Q) = 10 - Q$ and marginal cost 2: (a) the Stackelberg equilibrium with leader~1 and
follower~2; (b) the equilibrium when a third firm moves after firm~2, and how firms~1 and~2's
outputs change.

\method Backward induction with continuous actions: solve the last mover's best reply, substitute
it into the previous mover's profit, repeat (Tadelis \S8.3.2, eqs.~(8.1)--(8.2), pp.~160--161;
Gibbons \S2.1.B, pp.~61--64). The SPE lists each follower's whole reaction function, the outcome
only the quantities (Gibbons pp.~125--126).

\solution
\textbf{Notation.} $q_i \ge 0$ is firm $i$'s output, $Q$ the total; price-cost margin
$p - 2 = 8 - Q$; profit $\pi_i = (8 - Q)\,q_i$.

\textbf{(a) Two firms.}

\emph{Step 1: follower.} Firm~2 maximizes $\pi_2 = (8 - q_1 - q_2)\,q_2$:
\begin{equation}
  \frac{\partial \pi_2}{\partial q_2} = 8 - q_1 - 2q_2 = 0
  \;\Rightarrow\; q_2(q_1) = 4 - \frac{q_1}{2}\quad (q_1 \le 8). \label{eq:17r2}
\end{equation}
\emph{Step 2: leader.} Substituting \eqref{eq:17r2} into $\pi_1$:
\[
  \pi_1 = \Bigl(8 - q_1 - \bigl(4 - \tfrac{q_1}{2}\bigr)\Bigr) q_1
        = \Bigl(4 - \frac{q_1}{2}\Bigr) q_1 ,
  \qquad \frac{d\pi_1}{dq_1} = 4 - q_1 = 0 \;\Rightarrow\; q_1^* = 4 .
\]
\emph{Step 3: outcome.} From \eqref{eq:17r2}, $q_2^* = 4 - \frac{(4)}{2} = 2$; $Q = 6$,
$p = 10 - 6 = 4$, and
\[
  \pi_1 = (4 - 2)(4) = 8, \qquad \pi_2 = (4 - 2)(2) = 4 .
\]
SPE: $\bigl(q_1^* = 4,\ q_2(q_1) = \max\{0, 4 - q_1/2\}\bigr)$.

\textbf{(b) Three firms, moving in the order 1, 2, 3.}

\emph{Step 1: firm 3.} It maximizes $(8 - q_1 - q_2 - q_3)\,q_3$:
\begin{equation}
  8 - q_1 - q_2 - 2q_3 = 0 \;\Rightarrow\; q_3(q_1,q_2) = 4 - \frac{q_1 + q_2}{2}. \label{eq:17r3}
\end{equation}
\emph{Step 2: firm 2.} Substituting \eqref{eq:17r3} into $\pi_2$:
\[
  \pi_2 = \Bigl(8 - q_1 - q_2 - \bigl(4 - \tfrac{q_1+q_2}{2}\bigr)\Bigr) q_2
        = \Bigl(4 - \frac{q_1 + q_2}{2}\Bigr) q_2 ,
\]
\begin{equation}
  \frac{\partial \pi_2}{\partial q_2} = 4 - \frac{q_1}{2} - q_2 = 0
  \;\Rightarrow\; q_2(q_1) = 4 - \frac{q_1}{2}. \label{eq:17r2b}
\end{equation}
Substituting \eqref{eq:17r2b} into \eqref{eq:17r3}:
\[
  q_3(q_1) = 4 - \frac{q_1 + \bigl(4 - \frac{q_1}{2}\bigr)}{2} = 2 - \frac{q_1}{4}.
\]
\emph{Step 3: firm 1.} Substituting both:
\[
  \pi_1 = \Bigl(8 - q_1 - \bigl(4 - \tfrac{q_1}{2}\bigr) - \bigl(2 - \tfrac{q_1}{4}\bigr)\Bigr) q_1
        = \Bigl(2 - \frac{q_1}{4}\Bigr) q_1 ,
  \qquad \frac{d\pi_1}{dq_1} = 2 - \frac{q_1}{2} = 0 \;\Rightarrow\; q_1^* = 4 .
\]
\emph{Step 4: outcome.} $q_2^* = 4 - \frac{(4)}{2} = 2$, $q_3^* = 2 - \frac{(4)}{4} = 1$;
$Q = 7$, $p = 3$, margin $1$:
\[
  \pi_1 = 4,\qquad \pi_2 = 2,\qquad \pi_3 = 1 .
\]
\emph{Comparison.} Firms~1 and~2 produce exactly what they produced in (a), 4 and 2; firm~3 adds 1.
Total output rises from 6 to 7, the price falls from 4 to 3, and the profits of firms~1 and~2 halve
(8 to 4, and 4 to 2).

\interp Why do the outputs not move? Firm~3 absorbs half of any change in $q_1 + q_2$
(slope $-\frac12$ in \eqref{eq:17r3}). After substituting it, firm~2's residual margin is
$4 - \frac{q_1 + q_2}{2}$, half of the margin it faced in (a): the same problem scaled by $\frac12$,
so the same maximizer. The same scaling repeats upstream, which gives the pattern
$q_k = (a - c)/2^k$ for the $k$-th mover with demand $a - Q$ and marginal cost $c$ (here
$a - c = 8$: $4, 2, 1$). The leader always produces the
monopoly output $(a-c)/2 = 4$. Entry at the end of the chain does not change earlier outputs, but
it lowers their profits through the price.

\verify \texttt{check\_B\_L03.py} redoes every first-order condition with sympy (reaction functions,
$q^* = (4,2)$ and $(4,2,1)$, prices 4 and 3, profits $(8,4)$ and $(4,2,1)$, concavity of the
leader's problem) and repeats both games numerically with the leader on a $0.01$ grid. \emph{Comparison
with \MG{}'s answer} (pp.~186--189): same quantities, $(4,2)$ and $(4,2,1)$, and the same conclusion
that firms~1 and~2 do not change output. The book reports no prices or profits. Minor slips in its
text: the interior condition of firm~2's first-order condition is written ``if $q_1 > 0$'' where
$q_2 > 0$ is meant, both in (a) and in (b).

\answer{(a) $q_1 = 4$, $q_2 = 4 - q_1/2 = 2$; $p = 4$; $\pi = (8, 4)$. (b) $q_1 = 4$,
$q_2 = 4 - q_1/2 = 2$, $q_3 = 4 - (q_1+q_2)/2 = 1$; $p = 3$; $\pi = (4,2,1)$. Firms~1 and~2's outputs
are unchanged by firm~3's entry; their profits halve because the price falls.}
\end{exsol}
